\documentclass[12pt,oneside]{book}
\usepackage{../latex_compartido/estilo}
\renewcommand{\fechaEntregaActividad}{Por definir}
\begin{document}
\portadaActividad{14}{Paso de mensajes con UDP y TCP}
\frontmatter
\tableofcontents
\listoffigures
\listoftables
\mainmatter

\chapter{Análisis del paso de mensajes}
\textbf{Instrucción.} Identificar emisor, receptor, mensaje, canal, operaciones, direccionamiento, bloqueo y tipo de comunicación de cada implementación.

\section{Marco del problema}
Los programas de esta actividad son \textbf{procesos independientes}: no utilizan variables compartidas para intercambiar información. Se comunican enviando bytes a través de \textit{sockets} de red sobre la interfaz local (\texttt{localhost}). Un mensaje tiene un emisor, un receptor y un canal que lo transporta. El cliente solicita una operación y el servidor envía una respuesta.

\section{Análisis de UDP}
\begin{table}[H]\centering\small
\caption{Elementos de la implementación UDP.}
\begin{tabularx}{\textwidth}{P{3.4cm} L}
\toprule
\textbf{Elemento} & \textbf{Identificación}\\\midrule
Emisor y receptor & El cliente envía la solicitud; el servidor la recibe y después responde al cliente.\\
Mensaje & Texto UTF-8; la respuesta es \texttt{Mensaje recibido: [mensaje]}.\\
Canal & Datagramas UDP mediante \texttt{DatagramSocket}, puerto 5000.\\
Envío y recepción & \texttt{send()} y \texttt{receive()} con \texttt{DatagramPacket}.\\
Direccionamiento & Directo: \texttt{localhost:5000}. La respuesta se envía a la IP y puerto de origen del datagrama recibido.\\
Bloqueo & \texttt{receive()} espera un datagrama; la lectura del teclado espera entrada. \texttt{send()} no espera una confirmación de aplicación.\\
Comunicación & UDP no establece conexión ni garantiza entrega; \textit{nuestro diálogo} es síncrono, pues cada envío del cliente va seguido de una espera de respuesta.\\
\bottomrule
\end{tabularx}
\end{table}

\section{Análisis de TCP}
\begin{table}[H]\centering\small
\caption{Elementos de la implementación TCP.}
\begin{tabularx}{\textwidth}{P{3.4cm} L}
\toprule
\textbf{Elemento} & \textbf{Identificación}\\\midrule
Emisor y receptor & El cliente envía órdenes; el servidor responde con el resultado de cada orden.\\
Mensaje & Líneas de texto con \texttt{REGISTRAR}, \texttt{BUSCAR} o \texttt{SALIR}.\\
Canal & Conexión TCP creada con \texttt{Socket} y \texttt{ServerSocket}, puerto 6000.\\
Envío y recepción & \texttt{println()} escribe una línea; \texttt{readLine()} lee la siguiente.\\
Direccionamiento & Directo: \texttt{localhost:6000}; el sistema asigna un puerto temporal al cliente.\\
Bloqueo & \texttt{accept()} espera conexiones; \texttt{readLine()} espera una línea o cierre; conectar puede esperar o fallar.\\
Comunicación & Solicitud y respuesta síncronas al nivel de aplicación: el cliente espera el resultado antes de enviar otra orden. TCP ofrece un flujo fiable y ordenado de bytes.\\
\bottomrule
\end{tabularx}
\end{table}

\chapter{Servidor y cliente UDP}
\textbf{Instrucción.} Crear un servidor en el puerto 5000 que reciba datagramas, identifique dirección y puerto de origen, confirme su recepción y continúe hasta recibir \texttt{SALIR}; el cliente lee mensajes del teclado, los envía a \texttt{localhost} y muestra las respuestas.

\section{Entradas y procesamiento}
El usuario introduce una línea de texto, por ejemplo \texttt{Hola servidor}. El \texttt{ClienteUDP} la convierte a UTF-8 y crea un \texttt{DatagramPacket} con destino al puerto 5000. \texttt{ServidorUDP} permanece en \texttt{receive()} hasta obtener un datagrama; convierte sus bytes de nuevo a texto, muestra la dirección IP y el puerto del cliente y responde \texttt{Mensaje recibido: Hola servidor}. El cliente ejecuta también \texttt{receive()} para esperar dicha respuesta. Al enviar \texttt{SALIR}, el servidor confirma el mensaje y termina.

\section{Salidas, control y conceptos}
Las estructuras de control principales son \texttt{while}, \texttt{if} y el cierre automático de sockets mediante \texttt{try-with-resources}. El cliente tiene un tiempo máximo de espera de tres segundos por respuesta para evitar una espera ilimitada cuando no existe servidor. Esta limitación no convierte a UDP en un protocolo fiable.

\section{Diagrama de intercambio UDP}
\begin{center}
\begin{tabular}{r c l}
\textbf{Cliente UDP} & & \textbf{Servidor UDP}\\[.4em]
\texttt{send(datagrama)} & $\longrightarrow$ & \texttt{receive()}\\[.4em]
\texttt{receive()} & $\longleftarrow$ & \texttt{send(confirmación)}\\[.4em]
\texttt{SALIR} & $\longrightarrow$ & Recibe y responde\\[.4em]
Finaliza & $\longleftarrow$ & Confirma y finaliza
\end{tabular}
\end{center}

\section{Evidencia de código y terminales}
\espacioevidencia{figures/udp_codigo.png}{Código Java del cliente y servidor UDP; captura real pendiente.}
\espacioevidencia{figures/udp_servidor.png}{Terminal del servidor UDP; captura real pendiente.}
\espacioevidencia{figures/udp_cliente.png}{Terminal del cliente UDP; captura real pendiente.}

\chapter{Directorio de servicios TCP}
\textbf{Instrucción.} Implementar un servidor TCP con \texttt{ServerSocket(6000)} y un directorio \texttt{HashMap} que admita \texttt{REGISTRAR nombre puerto}, \texttt{BUSCAR nombre} y \texttt{SALIR}; construir un cliente interactivo para enviar las solicitudes y mostrar las respuestas.

\section{Entradas y procesamiento}
El servidor abre el puerto 6000 y utiliza \texttt{accept()} para aceptar un cliente. Conserva los nombres y puertos de servicios en un \texttt{HashMap<String,Integer>}. Por cada línea del cliente, separa sus palabras y ejecuta la acción indicada. Se valida que los puertos registrados sean números entre 1 y 65535. Una petición de búsqueda de un nombre ausente devuelve \texttt{Servicio no encontrado}. La orden \texttt{SALIR} devuelve \texttt{Conexion cerrada}, cierra la conexión y termina el servidor.

\section{Ejemplo esperado y salida observada}
\begin{table}[H]\centering\small
\caption{Solicitudes del directorio y respuestas.}
\begin{tabularx}{\textwidth}{L L}
\toprule
\textbf{Solicitud} & \textbf{Respuesta}\\\midrule
\texttt{REGISTRAR calculadora 7000} & Servicio registrado\\
\texttt{BUSCAR calculadora} & Servicio encontrado en el puerto 7000\\
\texttt{BUSCAR inventario} & Servicio no encontrado\\
\texttt{SALIR} & Conexion cerrada\\
\bottomrule
\end{tabularx}
\end{table}

El programa utiliza una estructura \texttt{HashMap}, ciclos \texttt{while}, decisiones \texttt{if/else}, lectura por líneas y un \texttt{try-with-resources}. Procesa una conexión a la vez, suficiente para el ejercicio. Los servicios se conservan en memoria mientras el servidor siga ejecutándose.

\section{Diagrama de intercambio TCP}
\begin{center}
\begin{tabular}{r c l}
\textbf{Cliente TCP} & & \textbf{Servidor TCP}\\[.4em]
\texttt{new Socket(...)} & $\longrightarrow$ & \texttt{accept()}\\[.4em]
\texttt{REGISTRAR ...} & $\longrightarrow$ & Guarda servicio\\[.4em]
Muestra respuesta & $\longleftarrow$ & \texttt{Servicio registrado}\\[.4em]
\texttt{BUSCAR ...} & $\longrightarrow$ & Consulta directorio\\[.4em]
Muestra respuesta & $\longleftarrow$ & Resultado de la búsqueda\\[.4em]
\texttt{SALIR} & $\longrightarrow$ & Confirma y cierra
\end{tabular}
\end{center}

\section{Evidencia de código y terminales}
\espacioevidencia{figures/tcp_codigo.png}{Código Java del cliente y servidor TCP; captura real pendiente.}
\espacioevidencia{figures/tcp_servidor.png}{Terminal del servidor TCP; captura real pendiente.}
\espacioevidencia{figures/tcp_cliente.png}{Terminal del cliente TCP; captura real pendiente.}

\chapter{Pruebas y comparación}
\section{Seis ejecuciones reproducibles}
Se ejecutaron \textbf{tres sesiones completas de cada implementación}, iniciando un proceso servidor y otro cliente en cada prueba. Los resultados siguientes se obtuvieron con el script \texttt{experiment001/run\_tests.py}, que almacena las salidas reales de ambas terminales en archivos de texto y verifica automáticamente las respuestas. La tabla no sustituye las capturas que deben tomarse en terminales separadas.

\begin{table}[H]\centering\scriptsize
\caption{Resultados de las seis sesiones cliente-servidor.}
\begin{tabularx}{\textwidth}{c c L L c}
\toprule
\textbf{Pr.} & \textbf{Proto.} & \textbf{Mensaje enviado} & \textbf{Respuesta} & \textbf{Resultado}\\\\midrule
1 & UDP & Hola servidor & Mensaje recibido: Hola servidor & Correcto\\
2 & UDP & Prueba de mensajes & Mensaje recibido: Prueba de mensajes & Correcto\\
3 & UDP & UDP funciona & Mensaje recibido: UDP funciona & Correcto\\
4 & TCP & BUSCAR calculadora & Servicio encontrado en el puerto 7000 & Correcto\\
5 & TCP & BUSCAR inventario & Servicio encontrado en el puerto 7001 & Correcto\\
6 & TCP & BUSCAR desconocido & Servicio no encontrado & Correcto\\
\bottomrule
\end{tabularx}
\end{table}

En cada sesión se envió también \texttt{SALIR} para cerrar el servidor. En las sesiones TCP de búsqueda exitosa, primero se registró el servicio correspondiente.

\section{Preguntas de análisis}
\textbf{1. ¿Qué métodos bloquean mientras esperan datos o conexiones?} En UDP, \texttt{DatagramSocket.receive()}. En TCP, \texttt{ServerSocket.accept()} y \texttt{BufferedReader.readLine()}; la apertura de la conexión puede bloquear durante el establecimiento. Leer del teclado también puede esperar entrada. Los métodos de envío no esperan una respuesta de la aplicación por sí mismos.

\textbf{2. ¿Por qué UDP no necesita \texttt{accept()}?} Porque no crea una conexión previa entre cliente y servidor. Cada datagrama contiene los datos necesarios para que el receptor identifique la dirección y el puerto de origen. En cambio, TCP establece una conexión y el servidor la recibe con \texttt{accept()}.

\textbf{3. ¿Qué diferencias existen entre \texttt{DatagramSocket}, \texttt{Socket} y \texttt{ServerSocket}?} \texttt{DatagramSocket} envía y recibe datagramas UDP independientes; \texttt{Socket} representa un extremo de una conexión TCP; y \texttt{ServerSocket} escucha un puerto y acepta nuevas conexiones TCP. TCP conserva el orden de los bytes recibidos y requiere delimitar los mensajes, aquí con saltos de línea.

\textbf{4. ¿Qué protocolo usarían para el directorio y por qué?} TCP, porque permite intercambiar varias órdenes por una conexión y proporciona entrega fiable y ordenada. Esto simplifica la implementación de un directorio que necesita recibir y contestar solicitudes correctamente. UDP requiere implementar controles adicionales si se desea garantizar entrega y reintentos.

\textbf{5. ¿Qué sucede cuando se inicia el cliente antes que el servidor?} En TCP, el intento de conexión suele fallar con \texttt{ConnectException: Connection refused} si no hay un proceso escuchando. En UDP, el envío puede finalizar sin error porque no hay establecimiento de conexión; el cliente espera una confirmación que no llegará y, en esta implementación, termina con \texttt{SocketTimeoutException} tras tres segundos.

\section{Conclusión}
El paso de mensajes permite colaborar a procesos independientes sin utilizar memoria compartida. UDP utiliza datagramas y no establece conexiones; TCP utiliza un flujo de bytes asociado a una conexión. En ambos ejercicios se programó un intercambio de solicitud y respuesta, con operaciones bloqueantes de recepción. Las tres pruebas UDP y las tres pruebas TCP completaron correctamente sus solicitudes. Los registros se conservan para reproducirlas; antes de entregar deberán añadirse capturas auténticas de ambas terminales.
\end{document}
